\chapter{Quantum Dynamics}
\label{ch:dynamics}

In non-relativistic quantum mechanics, time is a parameter labelling how states evolve, not an operator like $\hat{x}$. Measured quantities are expectation values $\langle\Psi|\hat{A}|\Psi\rangle$, and their time dependence can be placed in the state (the \emph{Schr\"odinger picture}) or in the operators (the \emph{Heisenberg picture}).

% ============================================================
\section{Schr\"odinger Dynamics}
% ============================================================

In the Schr\"odinger picture the state traces out a path in $\mathcal{H}$, $|\Psi,t\rangle = \sum_nc_n(t)|e_n\rangle$, while operators are fixed. Preserving the norm requires unitary evolution:
\[
    |\Psi,t\rangle = U(t,t_0)|\Psi,t_0\rangle,
    \qquad
    U^\dagger U = I .
\]
$U$ is unique and independent of the state, $U(t_0,t_0) = I$, it composes, $U(t_2,t_0) = U(t_2,t_1)U(t_1,t_0)$, and Hermitian conjugation reverses time, $U(t_0,t_1) = U^\dagger(t_1,t_0)$.

Differentiating, $\partial_t|\Psi,t\rangle = (\partial_tU)U^\dagger|\Psi,t\rangle$. Define
\[
    \hat{H}(t) \defeq i\hbar\,\pdv{U(t,t_0)}{t}U^\dagger(t,t_0).
\]
Differentiating $UU^\dagger = I$ shows $\hat{H} = \hat{H}^\dagger$, and the composition law shows it is independent of $t_0$. This gives the Schr\"odinger equation.

\begin{formula}[Schr\"odinger equation]
\[
    i\hbar\pdv{t}|\Psi,t\rangle = \hat{H}(t)|\Psi,t\rangle,
    \qquad
    i\hbar\pdv{t}U(t,t_0) = \hat{H}(t)U(t,t_0).
\]
\end{formula}

Any Hermitian operator can generate time evolution; the question is which physical system it describes. Classically $\frac{d}{dt}Q(x,p) = \{Q,H\}_{\text{PB}}$; quantum mechanically, for $\partial_t\hat{Q} = 0$,
\[
    \frac{d}{dt}\langle\hat{Q}\rangle = \frac{i}{\hbar}\langle[\hat{H},\hat{Q}]\rangle .
\]
So a classical system is quantized by promoting $x,p\to\hat{x},\hat{p}$ in $H(x,p)$. No classical correspondence is needed, though: for present-day quantum hardware the question is rather which Hamiltonians can be implemented reliably.

\begin{formula}[Solving for $U$]
\begin{enumerate}
    \item \textbf{Time-independent $\hat{H}$} (e.g.\ the oscillator): $U(t,t_0) = \exp\bigl[-\frac{i}{\hbar}\hat{H}(t-t_0)\bigr]$.
    \item \textbf{$[\hat{H}(t_1),\hat{H}(t_2)] = 0$} (e.g.\ spin in $\vect{B} = B(t)\khat$): $U(t,t_0) = \exp\bigl[-\frac{i}{\hbar}\int_{t_0}^tdt'\,\hat{H}(t')\bigr]$.
    \item \textbf{General case} (e.g.\ $\vect{B} = B_x(t)\ihat + B_z(t)\khat$): the time-ordered exponential
    \[
        U(t,t_0) = \mathcal{T}\exp\Bigl[-\frac{i}{\hbar}\int_{t_0}^tdt'\,\hat{H}(t')\Bigr]
        = 1 + \Bigl(-\frac{i}{\hbar}\Bigr)\int_{t_0}^tdt_1\hat{H}(t_1) + \Bigl(-\frac{i}{\hbar}\Bigr)^2\int_{t_0}^tdt_1\hat{H}(t_1)\int_{t_0}^{t_1}dt_2\hat{H}(t_2) + \cdots,
    \]
    ordered $t\ge t_1\ge t_2\ge\cdots\ge t_0$ (with no $1/n!$, since the ordering removes the permutations).
\end{enumerate}
\end{formula}
Case 2 works because $G = -\frac{i}{\hbar}\int\hat{H}$ commutes with $G' = -\frac{i}{\hbar}\hat{H}(t)$, so $\partial_te^G = G'e^G$ term by term. In case 3 this fails.

\begin{example}[Time evolution with a time-independent Hamiltonian]
If $\hat{H}|n\rangle = E_n|n\rangle$, then $U(t,0) = e^{-i\hat{H}t/\hbar} = \sum_ne^{-iE_nt/\hbar}|n\rangle\langle n|$ by the spectral theorem, and
\[
    |\Psi,t\rangle = \sum_ne^{-iE_nt/\hbar}|n\rangle\langle n|\Psi,0\rangle,
\]
recovering the expansion-in-eigenstates recipe of Chapter~\ref{ch:schrodinger}. In position space, $\langle x|U(t,0)|x'\rangle = K(x,t;x',0)$ is the \emph{propagator}: $\Psi(x,t) = \int dx'\,K(x,t;x',0)\Psi(x',0)$. For a free particle, inserting momentum states,
\[
    K = \int\frac{dp}{2\pi\hbar}e^{ip(x-x')/\hbar - ip^2t/2m\hbar} = \sqrt{\frac{m}{2\pi i\hbar t}}\exp\Bigl[\frac{im(x-x')^2}{2\hbar t}\Bigr].
\]
\end{example}

\begin{proof}[Why the time-ordered exponential works]
Differentiating the series term by term, $\partial_t$ acts only on the upper limit of the outermost integral, giving $-\frac{i}{\hbar}\hat{H}(t)$ times the series with one fewer integral, i.e.\ $-\frac{i}{\hbar}\hat{H}(t)U(t,t_0)$, with $\hat{H}(t)$ correctly on the left. Writing the $n$th term with all $n$ integrals running from $t_0$ to $t$ would overcount by $n!$, and require the integrand to be ordered, which is what $\mathcal{T}$ means.
\end{proof}

% ============================================================
\section{Heisenberg Dynamics}
% ============================================================

Since $\langle\Psi,t|\hat{A}_S|\Psi,t\rangle = \langle\Psi,t_0|U^\dagger\hat{A}_SU|\Psi,t_0\rangle$, we can instead evolve the operators.

\begin{definition}[Heisenberg operator]
\[
    \hat{A}_H(t) = U^\dagger(t,t_0)\,\hat{A}_S\,U(t,t_0).
\]
\end{definition}

Products and commutators carry over, $(\hat{A}\hat{B})_H = \hat{A}_H\hat{B}_H$ and $[\hat{A}_H,\hat{B}_H] = U^\dagger[\hat{A}_S,\hat{B}_S]U$; in particular $[\hat{x}_H(t),\hat{p}_H(t)] = i\hbar$. If $\hat{A}_S$ commutes with $\hat{H}_S$, then $\hat{A}_H = \hat{A}_S$; and $\hat{H}_H = \hat{H}_S$ when $\hat{H}$ commutes with itself at different times.

\begin{formula}[Heisenberg equation of motion]
\[
    i\hbar\frac{d}{dt}\hat{A}_H(t) = [\hat{A}_H(t),\hat{H}_H(t)] + i\hbar\Bigl(\pdv{\hat{A}_S}{t}\Bigr)_H .
\]
\end{formula}

\begin{theorem}[Ehrenfest's theorem]
For $\hat{H} = \hat{p}^2/2m + V(\hat{x})$,
\[
    \frac{d}{dt}\langle\hat{x}\rangle = \frac{\langle\hat{p}\rangle}{m},
    \qquad
    \frac{d}{dt}\langle\hat{p}\rangle = -\Bigl\langle\pdv{V}{x}\Bigr\rangle,
    \qquad
    m\frac{d^2}{dt^2}\langle\hat{x}\rangle = -\Bigl\langle\pdv{V}{x}\Bigr\rangle,
\]
the quantum analogue of Newton's second law. (Use $[\hat{x},\hat{p}^2] = 2i\hbar\hat{p}$ and $[\hat{p},V(\hat{x})] = -i\hbar V'(\hat{x})$.)
\end{theorem}

\begin{proof}[Derivation of the Heisenberg equation]
Differentiate $\hat{A}_H = U^\dagger\hat{A}_SU$, using $i\hbar\partial_tU = \hat{H}_SU$ and $-i\hbar\partial_tU^\dagger = U^\dagger\hat{H}_S$:
\[
    i\hbar\frac{d\hat{A}_H}{dt} = -U^\dagger\hat{H}_S\hat{A}_SU + U^\dagger\hat{A}_S\hat{H}_SU + i\hbar U^\dagger(\partial_t\hat{A}_S)U .
\]
Inserting $UU^\dagger = I$ between the operators turns the first two terms into $[\hat{A}_H,\hat{H}_H]$, with $\hat{H}_H = U^\dagger\hat{H}_SU$.
\end{proof}

Ehrenfest's theorem does not mean $\langle x\rangle$ follows a classical trajectory: in general $\langle V'(\hat{x})\rangle\neq V'(\langle x\rangle)$. They agree exactly for potentials at most quadratic (free particle, constant force, oscillator), and approximately when the wavepacket is narrow compared with the scale on which $V'$ varies.

\begin{example}[Free particle in the Heisenberg picture]
For $\hat{H} = \hat{p}^2/2m$, $\dot{\hat{p}}_H = 0$ and $\dot{\hat{x}}_H = \hat{p}_H/m$, so $\hat{x}_H(t) = \hat{x} + \hat{p}t/m$. Then
\[
    \Delta x^2(t) = \Delta x^2(0) + \frac{t}{m}\bigl(\langle\hat{x}\hat{p}+\hat{p}\hat{x}\rangle - 2\langle x\rangle\langle p\rangle\bigr) + \frac{t^2}{m^2}\Delta p^2,
\]
which for an initial Gaussian (no $x$--$p$ correlation) is $a^2 + \hbar^2t^2/4m^2a^2$, the spreading found before, now in two lines.
\end{example}

\begin{example}[Oscillator in the Heisenberg picture]
For the oscillator (Chapter~\ref{ch:sho}, with $|n\rangle$ the number states), $\frac{d}{dt}\hat{x}_H = \hat{p}_H/m$ and $\frac{d}{dt}\hat{p}_H = -m\omega^2\hat{x}_H$, so
\[
    \hat{x}_H(t) = \hat{x}\cos\omega t + \frac{\hat{p}}{m\omega}\sin\omega t,
    \qquad
    \hat{p}_H(t) = \hat{p}\cos\omega t - m\omega\hat{x}\sin\omega t,
\]
and $[\hat{H},\hat{a}] = -\hbar\omega\hat{a}$ gives $\hat{a}_H(t) = e^{-i\omega t}\hat{a}$, $\hat{a}^\dagger_H(t) = e^{i\omega t}\hat{a}^\dagger$. These look like the classical equations of motion, but in an energy eigenstate $\langle n|\hat{x}_H(t)|n\rangle = 0$: energy eigenstates do not behave classically.
\end{example}

\begin{example}[Spin precession in the Heisenberg picture]
For $\hat{H} = -\gamma B\hat{S}_z$, using $[\hat{S}_z,\hat{S}_x] = i\hbar\hat{S}_y$ and $[\hat{S}_z,\hat{S}_y] = -i\hbar\hat{S}_x$,
\[
    \dot{\hat{S}}_x = \frac{i}{\hbar}[\hat{H},\hat{S}_x] = \gamma B\hat{S}_y,
    \qquad
    \dot{\hat{S}}_y = -\gamma B\hat{S}_x,
    \qquad
    \dot{\hat{S}}_z = 0,
\]
so $\hat{S}_x(t) = \hat{S}_x\cos\omega t + \hat{S}_y\sin\omega t$ and $\hat{S}_y(t) = \hat{S}_y\cos\omega t - \hat{S}_x\sin\omega t$ with $\omega = \gamma B$. The operators themselves rotate, and any expectation value follows by taking $\langle\Psi,0|\cdot|\Psi,0\rangle$: the same Larmor precession found in the Schr\"odinger picture (Chapter~\ref{ch:twostate}), without ever writing down the state at time $t$.
\end{example}

% ============================================================
\section{The Interaction Picture}
% ============================================================

For $\hat{H} = \hat{H}_0 + \hat{V}(t)$ with $\hat{H}_0$ solvable, a third picture is useful: let $\hat{H}_0$ evolve the operators and $\hat{V}$ evolve the states. Define
\[
    |\Psi_I,t\rangle = e^{i\hat{H}_0t/\hbar}|\Psi,t\rangle,
    \qquad
    \hat{V}_I(t) = e^{i\hat{H}_0t/\hbar}\hat{V}(t)e^{-i\hat{H}_0t/\hbar}.
\]
Differentiating, the $\hat{H}_0$ terms cancel:
\[
    i\hbar\frac{d}{dt}|\Psi_I,t\rangle = \hat{V}_I(t)|\Psi_I,t\rangle .
\]
If $\hat{V} = 0$ the interaction-picture state is constant, and all change comes from the interaction. In the basis of $\hat{H}_0$ eigenstates, $\langle k|\hat{V}_I|n\rangle = e^{i(E_k-E_n)t/\hbar}V_{kn}$, which is exactly how the oscillating phases $e^{\pm2i\omega_0t}$ arose in the maser and the rotating frame arose in NMR. The interaction picture is the natural starting point for time-dependent perturbation theory.

% ============================================================
\section{Coherent States}
% ============================================================

Coherent states behave classically in many ways, and describe lasers and much of quantum optics.

\begin{definition}[Translation operator]
\[
    \hat{T}_{x_0} = \exp\Bigl(-\frac{i}{\hbar}x_0\hat{p}\Bigr).
\]
Translations form a commutative group, $\hat{T}_{x_0}\hat{T}_{y_0} = \hat{T}_{x_0+y_0}$, and $\hat{T}_{-x_0} = \hat{T}_{x_0}^\dagger = \hat{T}_{x_0}^{-1}$ is unitary. They translate the position operator, $\hat{T}_{x_0}^\dagger\hat{x}\hat{T}_{x_0} = \hat{x} + x_0$, and the states, $\hat{T}_{x_0}|x_1\rangle = |x_1+x_0\rangle$, so $\langle x|\hat{T}_{x_0}|\psi\rangle = \psi(x-x_0)$.
\end{definition}

\begin{definition}[Coherent state]
The coherent (Glauber) state is the translated oscillator ground state,
\[
    |\tilde{x}_0\rangle = \hat{T}_{x_0}|0\rangle,
    \qquad
    \langle x|\tilde{x}_0\rangle = \psi_0(x-x_0).
\]
\end{definition}

It is normalized, with $\langle\hat{x}\rangle = x_0$, $\langle\hat{p}\rangle = 0$ and $\langle\hat{H}\rangle = \frac12\hbar\omega + \frac12m\omega^2x_0^2$: the zero-point energy plus the classical potential energy at $x_0$. In the Heisenberg picture,
\[
    \langle\hat{x}_H(t)\rangle = x_0\cos\omega t,
    \qquad
    \langle\hat{p}_H(t)\rangle = -m\omega x_0\sin\omega t,
\]
classical harmonic motion, and $\Delta x^2 = \hbar/2m\omega$, $\Delta p^2 = m\hbar\omega/2$ at all times. So $\Delta x\,\Delta p = \hbar/2$ for all $t$: the Gaussian oscillates in the potential without changing shape.

With $d = \sqrt{\hbar/m\omega}$, write $\hat{T}_{x_0} = \exp\bigl(\frac{x_0}{\sqrt2d}(\hat{a}^\dagger-\hat{a})\bigr)$ and use the Baker--Campbell--Hausdorff formula $e^{A+B} = e^Ae^Be^{-\frac12[A,B]}$ (valid when $[A,B]$ is a number):
\[
    |\tilde{x}_0\rangle = e^{-x_0^2/4d^2}\sum_{n=0}^\infty\frac{1}{\sqrt{n!}}\Bigl(\frac{x_0}{\sqrt2d}\Bigr)^n|n\rangle .
\]
The occupation probabilities are Poisson, $P_n = \lambda^ne^{-\lambda}/n!$ with $\lambda = x_0^2/2d^2$, so $\langle N\rangle = \lambda$, $\Delta N = \sqrt\lambda$, and $\Delta H = \hbar\omega\sqrt\lambda$. For $x_0\gg d$, $\hbar\omega\ll\Delta H\ll\langle H\rangle$: the energy is sharp relative to its size, and the state becomes classical (Bohr's correspondence principle).

\begin{formula}[General coherent states]
With the displacement operator $D(\alpha) = \exp(\alpha\hat{a}^\dagger - \alpha^*\hat{a})$, $\alpha\in\mathbb{C}$,
\[
    |\tilde\alpha\rangle = D(\alpha)|0\rangle = e^{-|\alpha|^2/2}e^{\alpha\hat{a}^\dagger}|0\rangle,
    \qquad
    \hat{a}|\tilde\alpha\rangle = \alpha|\tilde\alpha\rangle,
\]
\[
    \langle\hat{x}\rangle = \sqrt2d\,\mathrm{Re}\,\alpha,
    \qquad
    \langle\hat{p}\rangle = \frac{\sqrt2\hbar}{d}\mathrm{Im}\,\alpha,
    \qquad
    |\tilde\alpha,t\rangle = e^{-i\omega t/2}|\widetilde{\alpha e^{-i\omega t}}\rangle .
\]
Coherent states are eigenstates of $\hat{a}$, which can serve as their definition; complex eigenvalues are allowed since $\hat{a}$ is not Hermitian.
\end{formula}

\begin{fact}[Coherent states are an overcomplete basis]
They are not orthogonal, $|\langle\tilde\alpha|\tilde\beta\rangle|^2 = e^{-|\alpha-\beta|^2}$, but they are complete,
\[
    I = \int\frac{d^2\alpha}{\pi}|\tilde\alpha\rangle\langle\tilde\alpha|,
\]
with $d^2\alpha/\pi = d\langle x\rangle\,d\langle p\rangle/2\pi\hbar$ the phase-space volume. In phase space a coherent state is a blob of fixed size $\Delta x\Delta p\sim\hbar$ circling the origin at frequency $\omega$, with number--phase uncertainty $\Delta N\sim|\alpha|$, $\Delta\varphi\sim1/|\alpha|$. A number state instead fills an annulus, $|\langle\tilde\alpha|n\rangle|^2 = e^{-|\alpha|^2}|\alpha|^{2n}/n!$.
\end{fact}

\begin{figure}[ht]
\centering
\begin{tikzpicture}[>=Stealth,scale=1.1]
  \draw[->] (-3,0) -- (3.2,0) node[right]{$\langle x\rangle/\sqrt2d$};
  \draw[->] (0,-3) -- (0,3.2) node[above]{$\langle p\rangle d/\sqrt2\hbar$};
  \draw[dashed] (0,0) circle (2.2);
  \fill[blue!30] (2.2,0) circle (0.35); \draw[blue!70!black] (2.2,0) circle (0.35);
  \fill[blue!30] (0,-2.2) circle (0.35); \draw[blue!70!black] (0,-2.2) circle (0.35);
  \draw[->,thick] (1.9,0.9) arc (25:-60:2.1);
  \node at (2.6,-1.6) {\small $e^{-i\omega t}$};
  \fill[red!20] (0,0) circle (1.25); \fill[white] (0,0) circle (0.95);
  \node[red!70!black] at (-1.6,1.3) {\small $|n\rangle$};
  \fill[green!25] (-1.5,-1.5) ellipse (0.15 and 0.6); \draw[green!50!black] (-1.5,-1.5) ellipse (0.15 and 0.6);
  \node[green!50!black] at (-2.3,-2.2) {\small squeezed};
\end{tikzpicture}
\caption{Phase-space pictures. A coherent state is a round blob of fixed area $\sim\hbar$ circling the origin at $\omega$; a number state is an annulus with no definite phase; a squeezed state is an ellipse, with reduced uncertainty in one quadrature at the expense of the other.}
\end{figure}

\begin{proof}[The coherent-state expansion via BCH]
With $\hat{T}_{x_0} = \exp\bigl(\frac{x_0}{\sqrt2d}(\hat{a}^\dagger-\hat{a})\bigr)$ and $A = \frac{x_0}{\sqrt2d}\hat{a}^\dagger$, $B = -\frac{x_0}{\sqrt2d}\hat{a}$, the commutator $[A,B] = \frac{x_0^2}{2d^2}$ is a number, so $e^{A+B} = e^Ae^Be^{-x_0^2/4d^2}$. Acting on $|0\rangle$, $e^B|0\rangle = |0\rangle$ since $\hat{a}|0\rangle = 0$, and $e^A|0\rangle = \sum_n\frac{1}{n!}\bigl(\frac{x_0}{\sqrt2d}\bigr)^n(\hat{a}^\dagger)^n|0\rangle = \sum_n\frac{1}{\sqrt{n!}}\bigl(\frac{x_0}{\sqrt2d}\bigr)^n|n\rangle$.
\end{proof}

\begin{figure}[ht]
\centering
\begin{tikzpicture}
\begin{axis}[width=10cm,height=5cm,xlabel={$t$},ylabel={$\langle x\rangle$},axis lines=middle,xmin=0,xmax=13,ymin=-1.5,ymax=1.5,ytick=\empty,xtick=\empty,samples=300,domain=0:13]
\addplot[name path=top,draw=none]{cos(deg(x))+0.25};
\addplot[name path=bot,draw=none]{cos(deg(x))-0.25};
\addplot[blue!20] fill between[of=top and bot];
\addplot[thick,blue!70!black]{cos(deg(x))};
\end{axis}
\end{tikzpicture}
\caption{A coherent state oscillates like a classical particle: $\langle x\rangle = x_0\cos\omega t$, with a constant width $\Delta x = \sqrt{\hbar/2m\omega}$ (shaded band).}
\end{figure}

% ============================================================
\section{Squeezed States}
% ============================================================

Another class of minimum-uncertainty states, \emph{squeezed states}, matters in quantum optics and is critical to Advanced LIGO. Suddenly change an oscillator from $(m_1,\omega_1)$ to $(m_2,\omega_2)$. The old ground state $|0\rangle_{(1)}$ is no longer stationary but still has $\Delta x\Delta p = \hbar/2$, with $\Delta x = e^{-\gamma}\sqrt{\hbar/2m_2\omega_2}$ and $e^\gamma = \sqrt{m_1\omega_1/m_2\omega_2}$. The two sets of ladder operators are related by a Bogoliubov transformation,
\[
    \begin{pmatrix}\hat{a}_1\\\hat{a}_1^\dagger\end{pmatrix} = \begin{pmatrix}\cosh\gamma&\sinh\gamma\\\sinh\gamma&\cosh\gamma\end{pmatrix}\begin{pmatrix}\hat{a}_2\\\hat{a}_2^\dagger\end{pmatrix}.
\]
Solving $\hat{a}_1|0\rangle_{(1)} = 0$ with the ansatz $|0\rangle_{(1)} = N\exp(-\frac12f\hat{a}_2^{\dagger2})|0\rangle_{(2)}$ gives $f = \tanh\gamma$ and $N = 1/\sqrt{\cosh\gamma}$.

\begin{formula}[Squeezed states]
For any oscillator,
\[
    |0\rangle_\gamma = \frac{1}{\sqrt{\cosh\gamma}}\exp\Bigl(-\frac12\tanh\gamma\,\hat{a}^{\dagger2}\Bigr)|0\rangle = S(\gamma)|0\rangle,
    \qquad
    S(\gamma) = \exp\Bigl(-\frac12(\gamma^*\hat{a}^2 - \gamma\hat{a}^{\dagger2})\Bigr),
\]
with $S(\gamma)$ unitary. More generally $|\tilde\alpha\rangle_\gamma = D(\alpha)S(\gamma)|0\rangle$, an eigenstate of $\hat{a}\cosh\gamma + \hat{a}^\dagger\sinh\gamma$ with eigenvalue $\alpha$. All are minimum-uncertainty states.
\end{formula}

As $\gamma\to\infty$, $|0\rangle_\infty\propto e^{-\hat{a}^{\dagger2}/2}|0\rangle$ is annihilated by $\hat{x}\propto\hat{a}+\hat{a}^\dagger$: it is the position eigenstate $|x = 0\rangle$, and translations give the others. As $\gamma\to-\infty$ one obtains momentum eigenstates. In phase space, squeezing deforms the round blob of the coherent state into an ellipse.

\paragraph{Quadrature variances.} Define the quadratures $\hat{X}_1 = \frac12(\hat{a}+\hat{a}^\dagger)$ and $\hat{X}_2 = \frac{1}{2i}(\hat{a}-\hat{a}^\dagger)$, dimensionless versions of $\hat{x}$ and $\hat{p}$ with $[\hat{X}_1,\hat{X}_2] = \frac i2$, so $\Delta X_1\Delta X_2\ge\frac14$. For a coherent state $\Delta X_1 = \Delta X_2 = \frac12$. The squeezed vacuum is annihilated by $\hat{b} = \hat{a}\cosh\gamma + \hat{a}^\dagger\sinh\gamma$, which satisfies $[\hat{b},\hat{b}^\dagger] = 1$; inverting, $\hat{a} = \hat{b}\cosh\gamma - \hat{b}^\dagger\sinh\gamma$. Then $\hat{X}_1 = e^{-\gamma}\cdot\frac12(\hat{b}+\hat{b}^\dagger)$ and $\hat{X}_2 = e^{\gamma}\cdot\frac{1}{2i}(\hat{b}-\hat{b}^\dagger)$, and since the state is the vacuum of $\hat{b}$,
\[
    \Delta X_1 = \tfrac12e^{-\gamma},
    \qquad
    \Delta X_2 = \tfrac12e^{\gamma}:
\]
one quadrature is squeezed below the vacuum noise and the other stretched, keeping the product at its minimum. Advanced LIGO injects squeezed vacuum into its interferometer to reduce the phase noise in the quadrature it measures.

% ============================================================
\section{The Quantized Electromagnetic Field}
% ============================================================

A single mode of the field in a cubic cavity of volume $V = L^3$ has
\[
    E_x(z,t) = \sqrt{\frac{2\omega^2}{\varepsilon_0V}}\,q(t)\sin kz,
    \qquad
    B_y(z,t) = \frac{\mu_0\varepsilon_0}{k}\sqrt{\frac{2\omega^2}{\varepsilon_0V}}\,p(t)\cos kz,
    \qquad
    k = \frac\omega c = \frac{n\pi}{L},
\]
where Maxwell's equations require $p = \dot{q}$. The energy is
\[
    H = \frac12\int dV\Bigl[\varepsilon_0E^2 + \frac{1}{\mu_0}B^2\Bigr] = \frac12\bigl(p^2 + \omega^2q^2\bigr),
\]
an oscillator with $m = 1$. Quantize by promoting $q,p$ to operators with $[\hat{q},\hat{p}] = i\hbar$, and define $\hat{a} = (\omega\hat{q} + i\hat{p})/\sqrt{2\hbar\omega}$, which removes a photon.

\paragraph{Why $p = \dot{q}$ and $H = \frac12(p^2+\omega^2q^2)$.} Faraday's law $\partial_zE_x = -\partial_tB_y$ with the mode functions above gives $-k\sqrt{\cdot}\,q\cos kz = -\frac{\mu_0\varepsilon_0}{k}\sqrt{\cdot}\,\dot{p}\cos kz$, so $\dot{p} = -\omega^2q$ (using $\omega = ck$), and the Amp\`ere--Maxwell law similarly gives $\dot{q} = p$. These are Hamilton's equations for $H = \frac12(p^2+\omega^2q^2)$, and the energy integral confirms it: $\int\sin^2kz\,dV = \int\cos^2kz\,dV = V/2$, so $\frac12\varepsilon_0\int E^2 = \frac12\omega^2q^2$ and $\frac{1}{2\mu_0}\int B^2 = \frac12p^2$.

\begin{formula}[Field operators]
\[
    \hat{E}_x(z,t) = \mathcal{E}_0\bigl(\hat{a}e^{-i\omega t} + \hat{a}^\dagger e^{i\omega t}\bigr)\sin kz,
    \qquad
    \mathcal{E}_0 = \sqrt{\frac{\hbar\omega}{\varepsilon_0V}} .
\]
\end{formula}

In a number state, $\langle n|\hat{E}_x|n\rangle = 0$ and $\Delta E_x = \mathcal{E}_0|\sin kz|\sqrt{2n+1}$: nothing like a classical field. (Since $\hat{a}$ and $\hat{a}^\dagger$ change $n$, their diagonal elements vanish; $\langle n|(\hat{a}e^{-i\omega t}+\hat{a}^\dagger e^{i\omega t})^2|n\rangle = \langle n|\hat{a}\hat{a}^\dagger + \hat{a}^\dagger\hat{a}|n\rangle = 2n+1$.) Even the vacuum $n = 0$ has $\Delta E_x\neq0$: these vacuum fluctuations drive spontaneous emission and the Lamb shift. In a coherent state,
\[
    \langle\tilde\alpha|\hat{E}_x|\tilde\alpha\rangle = 2\mathcal{E}_0|\alpha|\sin kz\cos(\omega t-\theta),
    \qquad
    \Delta E_x = \mathcal{E}_0|\sin kz|,
\]
with $\alpha = |\alpha|e^{i\theta}$, an oscillating field with fixed, $\alpha$-independent noise. For $|\alpha|\gg1$ it behaves like the classical field, and the single-mode field of a laser is idealized as a coherent state. Squeezed states trade noise between quadratures: amplitude-squeezed or phase-squeezed light. A general mode with wavevector $\lvect{k}$ and polarization $\lvect{e}_s$ is $\hat{\vect{E}}_{\lvect{k},s}\propto\lvect{e}_s\bigl(\hat{a}_{\lvect{k}}e^{i(\lvect{k}\cdot\lvect{r}-\omega t)} + \hat{a}_{\lvect{k}}^\dagger e^{-i(\lvect{k}\cdot\lvect{r}-\omega t)}\bigr)$, with $\omega = c|\lvect{k}|$.

\begin{fact}[Key results of this chapter]
\leavevmode
\begin{itemize}
\item $U(t,t_0)$ is unitary; $\hat{H} = i\hbar(\partial_tU)U^\dagger$; $U = e^{-i\hat{H}t/\hbar}$ for constant $\hat{H}$, time-ordered exponential in general.
\item Heisenberg picture: $i\hbar\dot{\hat{A}}_H = [\hat{A}_H,\hat{H}_H]$; Ehrenfest's theorem.
\item Coherent states $\hat{a}|\alpha\rangle = \alpha|\alpha\rangle$: Poisson number distribution, classical motion, minimum uncertainty at all times.
\item Squeezed states trade noise between quadratures; the EM field is a set of oscillators.
\end{itemize}
\end{fact}
